\documentclass[12pt,spanish,letterpaper,oneside]{article}
\def\documenttitle {Sistemas de protecci\'on y garant\'ias\\[6pt]de los derechos humanos}
\def\documentsubtitle {Articulaci\'on internacional, regional y constitucional}
\def\documenttitlehf {Sistemas de protecci\'on de derechos humanos}
\def\hfwidthwrap {true}
\def\hfwidthtitle {0.48}
\def\hfwidthcourse {0.46}
\def\documentsubject {Antecedentes de los derechos humanos}
\def\documentauthor {Martin Jonathan de la Cruz}
\def\coursename {Antecedentes de los derechos humanos}
\def\coursecode {005}
\def\universityname {Universidad Abierta y a Distancia de M\'exico}
\def\universityfaculty {Licenciatura en Derecho}
\def\universitydepartment {Antecedentes de los derechos humanos}
\def\universitydepartmentimage {departamentos/UnADM}
\def\universitydepartmentimagecfg {height=1.57cm}
\def\universitylocation {Ciudad de M\'exico}
\def\authortable {\begin{tabular}{ll}
\textbf{Alumno:} & Martin Jonathan de la Cruz \\
\textbf{Matr\'icula:} & ES2611202040 \\
\textbf{Figura docente:} & Daniel Alexis Lozano Ortega \\
\textbf{Semestre / Bloque:} & 2 / 2 \\
\textbf{Semana lectiva:} & Semana 7 \\
\textbf{Fecha de realizaci\'on:} & Noviembre de 2026 \\
\multicolumn{2}{l}{\universitylocation}
\end{tabular}}

\usepackage{iftex}
\ifLuaTeX
	\let\pdfpagewidth\pagewidth
	\let\pdfpageheight\pageheight
\fi

% Wrapper local para apuntar a la plantilla maestra de AulaTeX
\input{C:/Users/Sysx/Documents/AulaTeX-Academico/aulaTEX/base/Plantilla-Informe/template.tex}

\usepackage{helvet}
\renewcommand{\familydefault}{\sfdefault}
\usepackage{pdflscape}
\usepackage{array,tabularx,booktabs,longtable}
\usepackage{tikz}
\usetikzlibrary{shapes.geometric,arrows.meta,positioning,calc,fit,shadows,backgrounds,babel}
\usepackage{eso-pic}
\setcitestyle{authoryear,open={(},close={)}}

\newcommand{\insertcoverwatermark}{%
	\AddToShipoutPictureBG*{%
		\begin{tikzpicture}[remember picture,overlay]
			\node[opacity=0.16,inner sep=0pt] at (current page.center)
				{\includegraphics[width=0.72\paperwidth]{img/departamentos/UnADM.pdf}};
		\end{tikzpicture}%
	}%
}

\definecolor{azulOscuro}{HTML}{1A365D}
\definecolor{azulMedio}{HTML}{2B6CB0}
\definecolor{verdePetroleo}{HTML}{1A535C}
\definecolor{vinoOscuro}{HTML}{5C1D24}
\definecolor{grisFondo}{HTML}{F7FAFC}
\definecolor{bordeGris}{HTML}{CBD5E0}

\begin{document}
\insertcoverwatermark
\templatePortrait
\templatePagecfg
\onehalfspacing

\section{Introducci\'on}

Reconocer un derecho y garantizarlo son obligaciones distintas. La primera identifica una posici\'on jur\'idica de la persona; la segunda exige autoridades competentes, recursos efectivos y cumplimiento de las decisiones. Sin esas condiciones, el reconocimiento no impide por s\'i solo una detenci\'on arbitraria ni repara sus consecuencias. El art\'iculo 1 de la Constituci\'on mexicana obliga a prevenir, investigar, sancionar y reparar las violaciones de derechos humanos, pero sit\'ua esas obligaciones en el \`ambito de las competencias de cada autoridad y en los t\'erminos de la ley \citep[art. 1]{cpeumADH2026}.

Sostengo que la protecci\'on universal, interamericana y nacional debe evaluarse por su capacidad para limitar el poder, no por la cantidad de instituciones que la integran. Su complementariedad ampl\'ia las posibilidades de tutela, pero no elimina las diferencias entre una sentencia, una recomendaci\'on y un procedimiento de supervisi\'on. Tampoco convierte a los organismos internacionales en instancias de apelaci\'on ordinaria frente a cualquier resultado desfavorable.

La reforma constitucional de 2011 ampli\'o el par\'ametro de protecci\'on al reconocer los derechos contenidos en los tratados de los que M\'exico es parte. Sin embargo, la jurisprudencia P./J. 20/2014 incorpor\'o una salvedad decisiva: ante una restricci\'on constitucional expresa, debe atenderse al texto constitucional. Omitirla produce una descripci\'on incompleta del derecho mexicano; reconocerla no impide discutir cr\'iticamente sus efectos sobre la protecci\'on de la persona \citep{scjnParametro2014ADH}.

Ese debate adquiere especial importancia cuando el Estado castiga. La defensa de las v\'ictimas no habilita delitos imprecisos, penas retroactivas ni condenas sin prueba. Defiendo la \emph{ultima ratio} como criterio de contenci\'on: la intervenci\'on penal necesita una justificaci\'on de necesidad, adem\'as de fundamento legal. La legalidad penal, la presunci\'on de inocencia y la proporcionalidad no son concesiones al infractor; son l\'imites al poder estatal que protegen a cualquier persona \citep[arts. 14, 20 y 22]{cpeumADH2026}.

\clearpage
\section{Estructura org\'anica de los sistemas de protecci\'on}

Los tres sistemas comparten la protecci\'on de la persona, pero act\'uan mediante competencias y procedimientos diferentes. El Estado conserva la responsabilidad inmediata de hacer efectivos los derechos. La supervisi\'on internacional examina el cumplimiento de compromisos asumidos; no reemplaza la administraci\'on cotidiana de justicia ni absorbe las atribuciones de las instituciones nacionales.

\subsection{Fuentes, \`organos y alcance de la tutela}

En el sistema universal conviven mecanismos de naturaleza distinta. El Consejo de Derechos Humanos realiza el Examen Peri\'odico Universal y establece procedimientos especiales; el Alto Comisionado ofrece asistencia y seguimiento. Los comit\'es supervisan tratados espec\'ificos. La posibilidad de presentar una comunicaci\'on individual exige que el Estado haya aceptado la competencia correspondiente: ratificar un tratado no abre autom\'aticamente todos sus procedimientos de denuncia \citep{acnudhComunicaciones}.

En el sistema interamericano, las personas pueden presentar peticiones ante la Comisi\'on Interamericana de Derechos Humanos (CIDH). El acceso contencioso a la Corte Interamericana no es directo: corresponde a la Comisi\'on y a los Estados partes someter casos, bajo las condiciones convencionales. La Corte determina responsabilidad internacional del Estado y dispone reparaciones; no dicta condenas penales individuales \citep[arts. 44, 61--63 y 68]{oeaPacto1969ADH}.

En M\'exico, la tutela judicial incluye el amparo y los medios de control constitucional, con requisitos propios de procedencia y legitimaci\'on. La CNDH y los organismos locales tramitan quejas en su \`ambito de competencia y emiten recomendaciones p\'ublicas no vinculantes. Esta diferencia no vuelve irrelevante la v\'ia no jurisdiccional: permite documentar violaciones y exigir respuestas, pero no equivale a obtener una sentencia restitutoria \citep[arts. 102, apartado B, 103, 105 y 107]{cpeumADH2026}.

La figura~\ref{fig:sistemas} distingue fuentes, instituciones y resultados; la tabla~\ref{tab:matrizSistemas} precisa sus condiciones de acceso. Ninguna de las dos relaciones supone una jerarqu\'ia entre la ONU y la OEA.

\clearpage
\begin{landscape}
\thispagestyle{empty}
\begin{minipage}{\linewidth}
\singlespacing
\centering
\captionof{figure}{Sistemas de protecci\'on y garant\'ias de los derechos humanos}
\label{fig:sistemas}
\vspace{4pt}
\resizebox{\linewidth}{!}{%
\begin{tikzpicture}[
  x=1.0cm, y=1.0cm,
  headerTitle/.style={
    draw=azulOscuro, fill=azulOscuro, text=white, rounded corners=6pt,
    text width=24.5cm, align=center, inner sep=8pt,
    font=\small\bfseries
  },
  systemBox/.style={
    rounded corners=6pt, line width=1.3pt,
    inner sep=7pt
  },
  systemTitle/.style={
    rounded corners=4pt, text=white, font=\footnotesize\bfseries,
    inner sep=5pt, align=center
  },
  subBox/.style={
    rounded corners=4pt, fill=white, draw=bordeGris, line width=0.8pt,
    text width=7.1cm, align=left, inner sep=6pt, font=\fontsize{10}{12}\selectfont
  },
  normaBox/.style={
    rounded corners=4pt, fill=grisFondo, draw=bordeGris, line width=0.8pt,
    text width=7.1cm, align=left, inner sep=6pt, font=\fontsize{10}{12}\selectfont
  },
  connLine/.style={-Latex, line width=1.2pt}
]

% Titulo General
\node (title) at (12.0, 10.3) [headerTitle] {
  SISTEMAS DE PROTECCI\'ON Y GARANT\'IAS DE LOS DERECHOS HUMANOS\\[2pt]
  \scriptsize Articulaci\'on org\'anica entre el orden internacional universal, regional interamericano y constitucional nacional
};

% CAJA 1: SISTEMA UNIVERSAL (ONU)
\node (tUni) at (3.8, 8.7) [systemTitle, fill=azulMedio, text width=7.1cm] {
  1. SISTEMA UNIVERSAL (ONU)
};
\node (nUni) [normaBox, below=6mm of tUni] {
  \textbf{Instrumentos fundacionales:}\\
  $\bullet$ Carta de la ONU (1945) y Declaraci\'on Universal (DUDH, 1948).\\
  $\bullet$ Pactos de 1966: PIDCP y PIDESC; tratados tem\'aticos.
};
\node (oUni) [subBox, below=4mm of nUni] {
  \textbf{\'Organos principales de protecci\'on:}\\
  $\bullet$ \emph{Consejo de Derechos Humanos:} EPU y procedimientos especiales.\\
  $\bullet$ \emph{ACNUDH:} Asistencia t\'ecnica, observaci\'on e informes.
};
\node (mUni) [subBox, below=4mm of oUni] {
  \textbf{Comit\'es de tratados:}\\
  $\bullet$ Derechos Humanos, DESC, CEDAW y CAT, entre otros.\\
  $\bullet$ Informes y observaciones; comunicaciones individuales si el Estado acepta la competencia respectiva.\\
  $\bullet$ Dict\'amenes y seguimiento, no sentencias de una corte mundial.
};
\begin{scope}[on background layer]
\node (boxUni) [systemBox, draw=azulMedio, fill=azulMedio!5, fit=(tUni)(nUni)(oUni)(mUni)] {};
\end{scope}

% CAJA 2: SISTEMA REGIONAL INTERAMERICANO (OEA)
\node (tReg) at (12.0, 8.7) [systemTitle, fill=vinoOscuro, text width=7.1cm] {
  2. SISTEMA INTERAMERICANO (OEA)
};
\node (nReg) [normaBox, below=6mm of tReg] {
  \textbf{Par\'ametro convencional regional:}\\
  $\bullet$ Convenci\'on Americana (1969) y Protocolo de San Salvador (1988).\\
  $\bullet$ Declaraci\'on Americana de los Derechos y Deberes (1948).
};
\node (cReg) [subBox, below=4mm of nReg] {
  \textbf{Comisi\'on Interamericana (CIDH - Washington):}\\
  $\bullet$ Tramitaci\'on y examen de peticiones individuales.\\
  $\bullet$ Medidas cautelares urgentes e informes tem\'aticos y de pa\'is.\\
  $\bullet$ Soluci\'on amistosa y eventual sometimiento de casos a la Corte IDH.
};
\node (jReg) [subBox, below=4mm of cReg] {
  \textbf{Corte Interamericana (Corte IDH - San Jos\'e):}\\
  $\bullet$ Responsabilidad estatal, sentencias y reparaciones (arts. 63 y 68).\\
  $\bullet$ Opiniones consultivas y medidas provisionales.\\
  $\bullet$ Las personas no someten directamente casos contenciosos (art. 61).
};
\begin{scope}[on background layer]
\node (boxReg) [systemBox, draw=vinoOscuro, fill=vinoOscuro!5, fit=(tReg)(nReg)(cReg)(jReg)] {};
\end{scope}

% CAJA 3: SISTEMA CONSTITUCIONAL MEXICANO
\node (tNac) at (20.2, 8.7) [systemTitle, fill=verdePetroleo, text width=7.1cm] {
  3. SISTEMA NACIONAL (M\'EXICO)
};
\node (nNac) [normaBox, below=6mm of tNac] {
  \textbf{Par\'ametro de regularidad:}\\
  $\bullet$ Derechos constitucionales y de tratados; interpretaci\'on \emph{pro persona}.\\
  $\bullet$ Restricciones constitucionales expresas: criterio P./J. 20/2014.
};
\node (jNac) [subBox, below=4mm of nNac] {
  \textbf{Garant\'ias jurisdiccionales:}\\
  $\bullet$ Amparo directo e indirecto.\\
  $\bullet$ Acciones de inconstitucionalidad y controversias constitucionales.\\
  $\bullet$ Control por jueces federales y locales dentro de sus competencias.
};
\node (oNac) [subBox, below=4mm of jNac] {
  \textbf{Garant\'ias no jurisdiccionales:}\\
  $\bullet$ \emph{CNDH y organismos estatales} (Art. 102, Apartado B CPEUM).\\
  $\bullet$ Quejas por actos u omisiones de naturaleza administrativa p\'ublica.\\
  $\bullet$ Recomendaciones p\'ublicas no vinculantes. No sustituyen al amparo ni anulan sentencias.
};
\begin{scope}[on background layer]
\node (boxNac) [systemBox, draw=verdePetroleo, fill=verdePetroleo!5, fit=(tNac)(nNac)(jNac)(oNac)] {};
\end{scope}

% CAJA DE ARTICULACI\'ON Y PRINCIPIOS DE ENLACE
\node[fit=(boxUni)(boxReg)(boxNac), inner sep=0pt] (systems) {};
\node (boxArtic) [below=6mm of systems,
  draw=azulOscuro, fill=grisFondo, rounded corners=6pt, line width=1.1pt,
  text width=24.1cm, align=left, inner sep=6pt, font=\fontsize{10}{12}\selectfont
] {
  \textbf{Ejes de articulaci\'on procesal y sustantiva entre sistemas:}\\[2pt]
  \textbf{Acceso:} Agotamiento de recursos internos adecuados y efectivos, con las excepciones aplicables; no es requisito general del EPU ni de toda medida urgente.\\
  \textbf{Interpretaci\'on:} Protecci\'on m\'as favorable, competencia y motivaci\'on; la jurisprudencia interamericana se aplica conforme al criterio P./J. 21/2014.\\
  \textbf{L\'imite com\'un:} Proteger derechos no autoriza a crear delitos por analog\'ia ni a presumir culpabilidad.
};

% Conectores entre cajas
\draw[connLine, color=azulMedio] (title.south -| tUni.north) -- (tUni.north);
\draw[connLine, color=vinoOscuro] (title.south) -- (tReg.north);
\draw[connLine, color=verdePetroleo] (title.south -| tNac.north) -- (tNac.north);

\end{tikzpicture}
}
\par\vspace{6pt}
{\footnotesize\raggedright\textbf{Fuente:} Elaboraci\'on propia a partir de la CPEUM (arts. 1, 14 y 102--107), la CADH (arts. 44--68), el PIDCP, el PIDESC, el procedimiento de comunicaciones del ACNUDH y las tesis P./J. 20/2014 y 21/2014.\par}
\end{minipage}
\end{landscape}

\begin{landscape}
\thispagestyle{empty}
\begin{minipage}{\linewidth}
\singlespacing
\footnotesize
\centering
\captionof{table}{Competencias, acceso y efectos de los sistemas de protecci\'on}
\label{tab:matrizSistemas}
\renewcommand{\arraystretch}{1.35}
\begin{tabularx}{\linewidth}{>{\bfseries\raggedright\arraybackslash}p{2.7cm} *{3}{>{\raggedright\arraybackslash}X}}
\toprule
Elemento anal\'itico & Sistema Universal (ONU) & Sistema Regional (OEA) & Sistema Nacional (M\'exico) \\
\midrule
Instrumentos rectores & Carta ONU, DUDH, PIDCP, PIDESC y tratados tem\'aticos. & Convenci\'on Americana sobre DDHH y Protocolo de San Salvador. & Art\'iculos 1.º, 102 (B), 103, 105 y 107 de la CPEUM. \\
\addlinespace
\'Organos de tutela & Consejo de Derechos Humanos, ACNUDH y comit\'es de tratados. & Comisi\'on Interamericana y Corte Interamericana. & Justicia federal y local; CNDH y organismos locales en sus competencias. \\
\addlinespace
V\'ias de acceso & Comunicaciones ante comit\'es competentes; aportaciones a mecanismos de supervisi\'on. El EPU no resuelve demandas individuales. & Petici\'on ante CIDH (art. 44); solo Comisi\'on y Estados partes someten casos a la Corte (art. 61). & Amparo y recursos judiciales; quejas ante organismos de derechos humanos. El control abstracto tiene legitimaci\'on restringida. \\
\addlinespace
Condiciones & Aceptaci\'on de competencia, requisitos del tratado y agotamiento cuando proceda; excepciones y reglas sobre duplicidad internacional. & Agotamiento o excepci\'on acreditada, plazo y ausencia de duplicidad, conforme a arts. 46--47. & Procedencia, legitimaci\'on, plazos y competencia seg\'un la v\'ia elegida. La queja no reemplaza los recursos judiciales. \\
\addlinespace
Efectos & Dict\'amenes, observaciones y recomendaciones con seguimiento. No tienen por ese solo origen el r\'egimen de ejecuci\'on de una sentencia nacional. & Recomendaciones de CIDH y sentencias de Corte IDH: efectos distintos. El Estado parte en el caso debe cumplir la sentencia (art. 68). & Resoluciones judiciales obligatorias seg\'un su alcance; recomendaciones no vinculantes de CNDH y organismos locales. \\
\addlinespace
Relaci\'on con M\'exico & Derechos de tratados integrados al par\'ametro; valor de cada interpretaci\'on examinado seg\'un su fuente y competencia. & Jurisprudencia m\'as favorable, con verificaci\'on de razones aplicables y armonizaci\'on: P./J. 21/2014. & Art. 1 y par\'ametro constitucional; salvedad de restricciones expresas en P./J. 20/2014. \\
\bottomrule
\end{tabularx}
\par\vspace{6pt}
{\raggedright\textbf{Fuente:} CPEUM, CADH, procedimiento de comunicaciones del ACNUDH y jurisprudencias P./J. 20/2014 y 21/2014.\par}
\end{minipage}
\end{landscape}

\subsection{Subsidiariedad y alcance de la interpretaci\'on}

El agotamiento de recursos internos busca que el Estado atienda primero la vulneraci\'on. No debe transformarse en una exigencia de promover recursos in\'utiles indefinidamente. El art\'iculo 46.2 de la CADH prev\'e excepciones cuando falta el debido proceso para proteger el derecho, se impide acceder a los recursos o agotarlos, o existe retardo injustificado. Invocar una excepci\'on exige justificarla con hechos; afirmar que la justicia es ineficaz no basta \citep[art. 46]{oeaPacto1969ADH}.

Tampoco corresponde extender esa regla a todo mecanismo internacional. El EPU no es un litigio individual y la protecci\'on urgente obedece a presupuestos propios. En las comunicaciones ante comit\'es deben examinarse, adem\'as, la competencia aceptada, la materia, el momento de los hechos y las reglas sobre examen simult\'aneo de un mismo asunto \citep{acnudhComunicaciones}.

La jurisprudencia P./J. 21/2014 dispone que los criterios de la Corte Interamericana son vinculantes para los jueces mexicanos cuando resultan m\'as favorables a la persona. Si M\'exico no fue parte en el caso, deben verificarse las razones que hacen aplicable el precedente; tambi\'en debe procurarse su armonizaci\'on con la jurisprudencia nacional. No autoriza a reunir citas favorables prescindiendo de sus supuestos \citep{scjnInteramericana2014ADH}.

La tensi\'on con las restricciones constitucionales expresas no desaparece por invocar el principio \emph{pro persona}. Mi objeci\'on a una restricci\'on debe explicar qu\'e derecho limita, con qu\'e intensidad y por qu\'e resulta incompatible con una protecci\'on suficiente. Confundir esa cr\'itica con una facultad irrestricta para dejar de aplicar cualquier norma sustituir\'ia el control jur\'idico por la preferencia del int\'erprete.

\subsection{Estricta legalidad y contenci\'on del poder punitivo}

La estricta legalidad impide que la autoridad improvise los presupuestos del castigo. El art\'iculo 14 constitucional proh\'ibe imponer penas por analog\'ia o mayor\'ia de raz\'on y exige una ley exactamente aplicable; el art\'iculo 16 exige competencia, fundamentaci\'on y motivaci\'on. El PIDCP protege la legalidad penal y la presunci\'on de inocencia. Por eso, una finalidad p\'ublica defendible no subsana un tipo impreciso ni una imputaci\'on carente de prueba \citep[arts. 14 y 16]{cpeumADH2026}\citep[arts. 14--15]{onuPactoDerechosCivilesPoliticos1966}.

La \emph{ultima ratio} y la subsidiariedad internacional responden a preguntas diferentes. La primera exige justificar por qu\'e recurrir al instrumento penal frente a alternativas menos lesivas; la segunda regula la relaci\'on entre la tutela interna y ciertos procedimientos internacionales. No es necesario agotar un procedimiento administrativo para que todo delito pueda perseguirse, ni el car\'acter subsidiario de una petici\'on internacional convierte a la Corte Interamericana en tribunal penal.

Defiendo que el legislador demuestre la necesidad de criminalizar y que el juez exija los presupuestos legales de cada restricci\'on. Cuando medidas civiles, administrativas o preventivas protegen suficientemente el bien jur\'idico, la expansi\'on penal necesita una justificaci\'on adicional. Esta postura no supone atribuir al juez una facultad gen\'erica para ignorar delitos vigentes: exige cuestionar la norma por las v\'ias competentes y distinguir el juicio de constitucionalidad de la decisi\'on sobre responsabilidad individual.

La contenci\'on tampoco ampara la impunidad. El deber de investigar y sancionar las violaciones de derechos humanos coexiste con los derechos de defensa; no permite sustituir una investigaci\'on seria por la anticipaci\'on del castigo. El PIDCP establece que la prisi\'on preventiva no debe ser la regla general. Su uso como respuesta autom\'atica a la presi\'on p\'ublica desatiende esa exigencia y debilita la presunci\'on de inocencia \citep[art. 9.3]{onuPactoDerechosCivilesPoliticos1966}.

\clearpage
\section{Conclusiones}

La protecci\'on de los derechos humanos depende de instituciones que act\'uen dentro de su competencia y de decisiones que produzcan efectos verificables. La coexistencia de sistemas ampl\'ia las posibilidades de tutela, pero no dispensa de identificar el recurso adecuado, acreditar sus presupuestos y exigir su cumplimiento. Una recomendaci\'on puede contribuir a la reparaci\'on; no debe ofrecerse como si tuviera los efectos de una sentencia.

La apertura del derecho mexicano a las fuentes internacionales impone una argumentaci\'on m\'as rigurosa, no una legalidad m\'as flexible a conveniencia del poder. Debe reconocerse la salvedad de las restricciones constitucionales expresas y discutirse su alcance sin ocultarla. Tambi\'en debe explicarse por qu\'e un precedente interamericano resulta aplicable y m\'as protector. El prestigio de una instituci\'on no sustituye la raz\'on jur\'idica que sostiene su criterio.

Asumo una posici\'on contraria al uso del castigo como demostraci\'on de autoridad. La gravedad de un problema no prueba que una nueva pena sea necesaria, y la indignaci\'on frente a un hecho no acredita la culpabilidad de una persona. Defender la estricta legalidad exige rechazar la analog\'ia perjudicial, las imputaciones sin prueba y la normalizaci\'on de medidas restrictivas que deber\'ian justificarse individualmente. La \emph{ultima ratio} obliga a discutir la necesidad del instrumento penal, no solo la severidad de su aplicaci\'on.

Esa defensa no enfrenta los derechos de las v\'ictimas con los de las personas imputadas. Exige investigar con diligencia, sancionar cuando la responsabilidad se acredita mediante un proceso v\'alido y reparar las violaciones por las v\'ias correspondientes. Ninguno de esos deberes se cumple convirtiendo las garant\'ias en obst\'aculos prescindibles. Un Estado que invoca los derechos humanos para eludir los l\'imites de su propio poder contradice la obligaci\'on que dice proteger.

\clearpage
\begin{thebibliography}{99}

\bibitem[C\'amara de Diputados(2026)]{cpeumADH2026}
C\'amara de Diputados del H. Congreso de la Uni\'on. (2026). \emph{Constituci\'on Pol\'itica de los Estados Unidos Mexicanos} (\'ultima reforma publicada en el DOF el 7 de octubre de 2026). \url{https://www.diputados.gob.mx/LeyesBiblio/pdf/CPEUM.pdf}

\bibitem[Naciones Unidas(1948)]{onuDeclaracion1948ADH}
Naciones Unidas. (1948, 10 de diciembre). \emph{Declaraci\'on Universal de Derechos Humanos}. Resoluci\'on 217 A (III). \url{https://www.un.org/es/about-us/universal-declaration-of-human-rights}

\bibitem[Naciones Unidas(1966a)]{onuPactoDerechosCivilesPoliticos1966}
Naciones Unidas. (1966a, 16 de diciembre). \emph{Pacto Internacional de Derechos Civiles y Pol\'iticos}. \url{https://www.ohchr.org/sites/default/files/ccpr_SP.pdf}

\bibitem[Naciones Unidas(1966b)]{onuPactoDerechosEconomicosSocialesCulturales1966}
Naciones Unidas. (1966b, 16 de diciembre). \emph{Pacto Internacional de Derechos Econ\'omicos, Sociales y Culturales}. \url{https://www.ohchr.org/sites/default/files/cescr_SP.pdf}

\bibitem[ACNUDH(s.\,f.)]{acnudhComunicaciones}
Oficina del Alto Comisionado de las Naciones Unidas para los Derechos Humanos. (s.\,f.). \emph{Procedimiento de comunicaciones individuales ante los \`organos de tratados}. \url{https://www.ohchr.org/es/treaty-bodies/individual-communications}

\bibitem[Organizaci\'on de los Estados Americanos(1969)]{oeaPacto1969ADH}
Organizaci\'on de los Estados Americanos. (1969, 22 de noviembre). \emph{Convenci\'on Americana sobre Derechos Humanos}. \url{https://www.oas.org/dil/esp/tratados_b-32_convencion_americana_sobre_derechos_humanos.htm}

\bibitem[SCJN(2014a)]{scjnParametro2014ADH}
Suprema Corte de Justicia de la Naci\'on. (2014a). \emph{Derechos humanos contenidos en la Constituci\'on y en los tratados internacionales. Constituyen el par\'ametro de control de regularidad constitucional, pero cuando en la Constituci\'on haya una restricci\'on expresa al ejercicio de aqu\'ellos, se debe estar a lo que establece el texto constitucional}. Tesis P./J. 20/2014 (10a.), registro 2006224. Gaceta del Semanario Judicial de la Federaci\'on, libro 5, tomo I, 202. \url{https://sjf2.scjn.gob.mx/detalle/tesis/2006224}

\bibitem[SCJN(2014b)]{scjnInteramericana2014ADH}
Suprema Corte de Justicia de la Naci\'on. (2014b). \emph{Jurisprudencia emitida por la Corte Interamericana de Derechos Humanos. Es vinculante para los jueces mexicanos siempre que sea m\'as favorable a la persona}. Tesis P./J. 21/2014 (10a.), registro 2006225. Gaceta del Semanario Judicial de la Federaci\'on, libro 5, tomo I, 204. \url{https://sjf2.scjn.gob.mx/detalle/tesis/2006225}

\end{thebibliography}

\end{document}
